\begin{env}[3.7.2]
\label{0.3.7.2-fr}
Conservando las notaciones y las hipótesis de (3.5.3), supongamos que
$\mathcal{G}$ es un \emph{haz} de grupos (resp. de anillos) sobre $Y$.
La definición de las secciones de $\psi^*(\mathcal{G})$ (3.7.1) muestra
(teniendo en cuenta (3.4.4)) que el homomorfismo de tallos
$\psi_x\circ\rho_{\psi(x)}:\mathcal{G}_{\psi(x)}\to
(\psi^*(\mathcal{G}))_x$ es un \emph{isomorfismo funtorial en
$\mathcal{G}$}, que permite identificar estos dos tallos. Con esta
identificación, $u_x^\sharp$ es idéntico al homomorfismo definido en
(3.5.1) y, en particular, se tiene
$\operatorname{Supp}(\psi^*(\mathcal{G}))=
\psi^{-1}(\operatorname{Supp}(\mathcal{G}))$.

Una consecuencia inmediata de este resultado es que \emph{el funtor
$\psi^*(\mathcal{G})$ es exacto en $\mathcal{G}$} en la categoría abeliana
de los haces de grupos conmutativos.
\end{env}
